\section{Limitations}\label{sec:limits}
\textbf{Data scale and label skew.} The benchmark has 53 patients; the prescribed split leaves 603 training
annotations and only 8 test patients (206 annotations). Rare labels are nearly absent: the test split has no
\emph{excluded} example and only 7 \emph{not included}, so per-class scores for these classes are either
undefined or rest on a handful of cases, and patient-level bootstrap intervals are wide.

\textbf{Single seed, limited search.} Each training configuration was run once (seed 42) and the
hyper-parameter search was deliberately small to respect the 6\,GB budget and the deadline; reported
differences of a few points are within run-to-run variance we did not measure.

\textbf{Explanations.} Rationales are constrained to the cited sentences and pass an automatic grounding check,
but they were not rated by clinicians, and the check (numbers and label contradiction) cannot catch every
unfaithful paraphrase. They should be read as summaries of the cited evidence, not as clinical reasoning.

\textbf{Inference measurements.} Timings are from one consumer GPU under Windows with a slow host CPU; at
batch size 1 small kernels are launch-bound, which shapes every speed result (Section~\ref{sec:results}). The
tiled attention kernel is an educational PyTorch implementation, not a fused kernel, and was validated in
isolation rather than substituted into the production model path.

\textbf{Aggregation and scope.} The overall outcome is a hand-written rule over criterion predictions; the
dataset contains no patient--trial ground truth, so this outcome is not evaluated. Criteria fetched from
ClinicalTrials.gov are split heuristically. \sys\ is an educational screening prototype and must not be used
to make enrollment or treatment decisions.
